% 352_Leptonreste_En_ch.tex
% Johann Pascher, 5. September 2026
\providecommand{\BM}{\textbf{[B]}}
\providecommand{\KM}{\textbf{[K]}}
\providecommand{\SM}{\textbf{[S]}}
\providecommand{\QM}{\textbf{[Q]}}

\phantomsection
\addcontentsline{toc}{chapter}{Lepton Residuals: Lepton Ladder, Galois and Koide Compared}

\section*{Status Markers}
\begin{center}
\begin{tabular}{@{}lp{10.5cm}@{}}
\textbf{[K]} & \emph{Confirmed} --- derived from $\xi$, numerically verified. \\[4pt]
\textbf{[B]} & \emph{Proved} --- algebraically established. \\[4pt]
\textbf{[S]} & \emph{Speculative} --- open. \\[4pt]
\textbf{[Q]} & \emph{Source} --- corpus-external primary source, verified here. \\
\end{tabular}
\end{center}

\section*{Abstract}

Three independent FFGFT structures make predictions for lepton ratios: the lepton ladder
(Docs.~006/317), the Galois identities (Doc.~351), and the Koide relation. This document
compares all three predictions with PDG values and expresses residuals in units of
$\xi = 4/30000$. All formulae are exactly rational; numerical evaluations use
$m_e = 0.51099895$~MeV, $m_\mu = 105.6583755$~MeV, $m_\tau = 1776.86$~MeV (PDG) \QM.

\subsection*{Theoretical Claim}

FFGFT claims predictions for dimensionless ratios, not for dimensional individual masses
(Doc.~351, §\,Theoretical Claim). Individual masses carry $K_{\rm frak}$, SI projection,
and QED-dependent extraction; these produce an unavoidable bandwidth of order $\xi$ to
$100\xi$. Ratios, by contrast, are correction-free --- \emph{provided the corrections
are uniform} (see §\,9).

\section{Mass Formula and Quantum Numbers}

The FFGFT mass formula reads (Docs.~006, 333):
\begin{equation}
  m_i = r_i \cdot \xi^{p_i} \cdot v \cdot K_{\rm frak} \qquad \KM
\end{equation}
with $\xi = 4/30000$ \BM, $K_{\rm frak} = 74/75$ \KM. In the ratio of two leptons
$v$ and $K_{\rm frak}$ cancel:
\begin{equation}
  \frac{m_j}{m_i} = \frac{r_j}{r_i} \cdot \xi^{p_j - p_i} \qquad \BM
\end{equation}

\begin{center}
\begin{tabular}{@{}lllll@{}}
\toprule
Particle & $r_i$ & $p_i$ & $n_\theta,\,n_\phi$ & $r_i = n_\phi^2/n_\theta$ \\
\midrule
Electron $e$ & $4/3$ & $3/2$ & $3,\,2$ & $4/3$ \\
Muon $\mu$   & $16/5$ & $1$ & $5,\,4$ & $16/5$ \\
Tau $\tau$   & $25/9$ & $2/3$ & $9,\,5$ & $25/9$ \\
\bottomrule
\end{tabular}
\end{center}

\section{Lepton Ladder Predictions}

\begin{equation}
  \frac{m_\mu}{m_e} = \frac{16/5}{4/3} \cdot \xi^{1-3/2}
  = \frac{12}{5} \cdot \xi^{-1/2}
  = \frac{12}{5} \cdot \sqrt{7500}
\end{equation}

\begin{center}
\begin{tabular}{@{}lllll@{}}
\toprule
Ratio & Exact formula & FFGFT & PDG & Residual \\
\midrule
$m_\mu/m_e$ & $\tfrac{12}{5}\cdot\xi^{-1/2}$ & $207.846$ & $206.768$ & $+39.1\cdot\xi$ \KM \\
$m_\tau/m_\mu$ & $\tfrac{125}{144}\cdot\xi^{-1/3}$ & $16.992$ & $16.817$ & $+77.9\cdot\xi$ \KM \\
$m_\tau/m_e$ & $\tfrac{25}{12}\cdot\xi^{-5/6}$ & $3531.6$ & $3477.2$ & $+117.4\cdot\xi$ \KM \\
\bottomrule
\end{tabular}
\end{center}

Algebraic dependence \BM:
\begin{equation}
  \varepsilon(\tau/e) \approx \varepsilon(\tau/\mu) + \varepsilon(\mu/e)
  = 77.9 + 39.1 = 117.0\cdot\xi
\end{equation}
(The exact value $117.4$ contains a cross-term of order $\xi^2$.)
There are only two independent residuals.

\section{Galois Identities}

From the Galois field structure (Doc.~351) \BM:
\begin{align}
  (m_\mu/m_e)^2 &= 43200 = |GF(9)^*|^2 \cdot 5^2 \cdot |GF(27)| \\
  m_e \cdot m_\mu &= 54~\text{MeV}^2 \quad(\text{Layer-2 with anchor }m_e)
\end{align}

\begin{center}
\begin{tabular}{@{}lllll@{}}
\toprule
Identity & Galois & PDG & Residual & Status \\
\midrule
$(m_\mu/m_e)^2$ & $43200$ & $42753.1$ & $-77.6\cdot\xi$ & \KM; cause \SM \\
$m_e\cdot m_\mu$~[MeV²] & $54$ & $53.991$ & $-1.21\cdot\xi$ & emp.\ discrepancy \KM \\
\bottomrule
\end{tabular}
\end{center}

\textbf{Connection to the lepton ladder.} The Galois difference residual $-77.6\cdot\xi$
equals twice the linear lepton ladder residual ($+39.1\cdot\xi$) with opposite sign ---
algebraically forced, since $(m_\mu/m_e)^2_{\rm FFGFT} > (m_\mu/m_e)^2_{\rm PDG}$
if and only if $m_\mu/m_e|_{\rm FFGFT} > m_\mu/m_e|_{\rm PDG}$.

\section{Koide Relation}

\begin{equation}
  Q \equiv \frac{m_e + m_\mu + m_\tau}{(\sqrt{m_e} + \sqrt{m_\mu} + \sqrt{m_\tau})^2}
  \stackrel{?}{=} \frac{2}{3}
\end{equation}

\begin{center}
\begin{tabular}{@{}llll@{}}
\toprule
Input & $Q$ & Deviation & Status \\
\midrule
PDG masses & $0.666\,660\,5$ & $-0.046\cdot\xi$ & \KM \\
FFGFT predictions (anchor $m_e$) & $0.667\,742$ & $+8.06\cdot\xi$ & numerical \\
\bottomrule
\end{tabular}
\end{center}

\textbf{Koide as $m_\tau$ prediction.} Given PDG $m_e$, $m_\mu$ and $Q = 2/3$,
algebraically \KM:
\begin{equation}
  m_{\tau,{\rm Koide}} = 1776.969~\text{MeV}
  \quad\text{vs.}\quad
  m_{\tau,{\rm PDG}} = 1776.860~\text{MeV}
  \qquad\text{deviation: }+0.46\cdot\xi
\end{equation}
This is the most precise prediction in the entire comparison.

\section{Overall Comparison}

\begin{center}
\begin{tabular}{@{}llllll@{}}
\toprule
Quantity & Method & Prediction & PDG & Residual & Status \\
\midrule
$m_\mu/m_e$ & Ladder & $207.846$ & $206.768$ & $+39.1\cdot\xi$ & \KM \\
$m_\tau/m_\mu$ & Ladder & $16.992$ & $16.817$ & $+77.9\cdot\xi$ & \KM \\
$m_\tau/m_e$ & Ladder & $3531.6$ & $3477.2$ & $+117.4\cdot\xi$ & \KM; sum \\
$(m_\mu/m_e)^2$ & Galois & $43200$ & $42753$ & $-77.6\cdot\xi$ & \KM; cause \SM \\
$m_e\cdot m_\mu$~[MeV²] & Galois & $54$ & $53.991$ & $-1.21\cdot\xi$ & emp.\ disc.\ \KM \\
$Q_{\rm Koide}$ & Koide & $2/3$ & $0.66666$ & $-0.046\cdot\xi$ & \KM \\
$m_\tau$~[MeV] & Koide & $1776.97$ & $1776.86$ & $+0.46\cdot\xi$ & \KM \\
\bottomrule
\end{tabular}
\end{center}

\section{Findings}

\textbf{(1) Koide is most precise \KM.}
PDG masses satisfy Koide to $0.046\cdot\xi$. Koide gives with $+0.46\cdot\xi$ the
most accurate $m_\tau$ prediction in the comparison. FFGFT predictions meet Koide
only to $+8\cdot\xi$.

\textbf{(2) Lepton ladder and Galois share the same basic residual \KM.}
The Galois difference residual ($-77.6\cdot\xi$) equals exactly $-2\times(+39.1\cdot\xi)$
of the lepton ladder residual for $m_\mu/m_e$. This is algebraically forced.

\textbf{(3) All residuals positive except Galois \KM.}
Lepton ladder and Koide-$m_\tau$ lie just above PDG; Galois-$(m_\mu/m_e)^2$ and
$Q_{\rm Koide}$ lie just below target. Signs are consistent.

\textbf{(4) Range of residuals \KM.}
\begin{center}
\begin{tabular}{@{}ll@{}}
\toprule
Method & Residual \\
\midrule
Koide ($Q$) & $\approx 0.05\cdot\xi$ \\
Galois (product) & $\approx 1\cdot\xi$ \\
Koide ($m_\tau$) & $\approx 0.5\cdot\xi$ \\
Lepton ladder ($m_\mu/m_e$) / Galois (square) & $\approx 39$--$78\cdot\xi$ \\
\bottomrule
\end{tabular}
\end{center}

\textbf{(5) Koide is scale-invariant \BM.}
For every $\lambda > 0$:
\begin{equation}
  Q(\lambda m_1, \lambda m_2, \lambda m_3) = Q(m_1, m_2, m_3) \quad\BM
\end{equation}
Consequence: no uniform factor --- neither $v$, nor $K_{\rm frak}$, nor a common
SI projection --- can change $Q$. Option~2 (corrections improve Koide) is thus
algebraically excluded.

\section{Scale Invariance and Its Consequences}

\subsection*{What Koide can measure}

$Q$ depends exclusively on the ratios of masses, not on their absolute value:
\begin{equation}
  Q = Q\!\left(1,\,\frac{m_\mu}{m_e},\,\frac{m_\tau}{m_e}\right)
\end{equation}

With the FFGFT ratios $(m_\mu/m_e)_{\rm FFGFT} = 207.85$ and
$(m_\tau/m_e)_{\rm FFGFT} = 3531.6$, $Q_{\rm FFGFT} = 0.66774$ is fully determined
--- independent of $v$, $K_{\rm frak}$, or the SI projection.

\subsection*{Three tested options}

\begin{center}
\begin{tabular}{@{}lll@{}}
\toprule
Option & $Q$ change & Status \\
\midrule
Uniform $K_{\rm frak}$ (all equal) & none & \BM\ excluded \\
Single $v$ anchor & none & \BM\ excluded \\
Generation-dep.\ $K^{\alpha p_i}_{\rm frak}$ & yes, but worse masses & fit without physics \\
\bottomrule
\end{tabular}
\end{center}

\textbf{Generation-dependent test.} The exponent $\alpha$ in $m_i\cdot K^{\alpha p_i}_{\rm frak}$
that forces $Q = 2/3$ lies at $\alpha\approx -36/17$. The resulting masses
($m_e = 0.527$~MeV, $m_\mu = 108.0$~MeV, $m_\tau = 1817.6$~MeV) deviate further
from PDG than the uncorrected bare masses. No physics behind this value.

\subsection*{Conclusion}

Koide is not contained in the FFGFT lepton ladder \SM. The lepton ladder organises
masses via logarithmic $\xi$ powers: $m_i \propto \xi^{p_i}$. Koide organises masses
via square roots $\sqrt{m_i}$. These two structures are orthogonal --- there is no
algebraic relation between $\xi^{p_i}$ and $\sqrt{r_i\,\xi^{p_i}}$ that forces Koide.

PDG masses satisfy Koide to $0.046\cdot\xi$ because the three PDG masses as a
measurement system are mutually consistent --- they carry implicitly a symmetry not
yet anchored in the FFGFT quantum numbers $(r_i, p_i)$.

The natural candidate for this anchoring is the $Z_3$ circulant structure of the lepton
mass operator (Doc.~A110): there $\theta = 2/9$ describes the Koide phase in the angular
sector.

\section{Koide Condition in the FFGFT Quantum Numbers}

\subsection*{$Z_3$ parametrisation}

The eigenvalues of a $Z_3$ circulant mass operator can be written as:
\begin{equation}
  \sqrt{m_k} = c + d\cdot\cos\!\left(\theta + \frac{2\pi k}{3}\right), \quad k = 0,1,2
\end{equation}
Then $\sum\sqrt{m_k} = 3c$ and $\sum m_k = 3c^2 + \tfrac{3}{2}d^2$, giving:
\begin{align}
  Q &= \frac{1}{3} + \frac{d^2}{6c^2} \\
  Q = \frac{2}{3} &\iff d = c\sqrt{2} \quad\BM
\end{align}

\subsection*{$d/c$ ratio in PDG and FFGFT}

\begin{center}
\begin{tabular}{@{}llll@{}}
\toprule
Source & $d/c$ & $d/c - \sqrt{2}$ & Status \\
\midrule
PDG & $1.41420$ & $-0.098\cdot\xi$ & \KM \\
FFGFT predictions & $1.41649$ & $+17.09\cdot\xi$ & numerical \\
Target & $\sqrt{2} = 1.41421$ & $0$ & --- \\
\bottomrule
\end{tabular}
\end{center}

PDG satisfies $d = c\sqrt{2}$ to $0.10\cdot\xi$. FFGFT predictions miss it by
$+17\cdot\xi$.

\subsection*{Angle $\theta$ and connection to Doc.~A110}

With assignment $k=0\to\tau$, $k=1\to e$, $k=2\to\mu$:

\begin{center}
\begin{tabular}{@{}lll@{}}
\toprule
Source & $\theta$~[rad] & $\theta$~[°] \\
\midrule
PDG & $0.22223$ & $12.733$ \\
FFGFT predictions & $0.22099$ & $12.662$ \\
Doc.~A110 ($\theta = 2/9$) & $0.22222$ & $12.732$ \\
\bottomrule
\end{tabular}
\end{center}

All three agree to $\lesssim 0.5\cdot\xi$ \KM. The Koide phase angle $\theta = 2/9$
from Doc.~A110 correctly describes the angular structure. The amplitude $d/c$,
however, is not determined by $\theta$ --- it is an independent condition.

\subsection*{Koide condition in the quantum numbers $(r_i, p_i)$}

Since $x_i = \sqrt{r_i}\cdot\xi^{p_i/2}$ and Koide is scale-invariant, $Q$ depends
only on the differences of exponents:
\begin{equation}
  \Delta p_e = p_e - p_\tau = \frac{5}{6},\quad
  \Delta p_\mu = p_\mu - p_\tau = \frac{1}{3},\quad
  \Delta p_\tau = 0
\end{equation}

The exact Koide condition $Q = 2/3$ requires:
\begin{equation}
  r_{\mu,{\rm Koide}} = 3.2378\ldots \quad\text{instead of}\quad r_\mu = \frac{16}{5} = 3.2000
\end{equation}
Deviation: $\Delta r_\mu / r_\mu = +88.6\cdot\xi$. The nearest rational approximation
is $531/164$ --- no representation as $n_\phi^2/n_\theta$ with small integer winding numbers.

\begin{center}
\begin{tabular}{@{}llll@{}}
\toprule
Particle & FFGFT $r_i$ & Winding $(n_\theta, n_\phi)$ & Koide-$r_i$ \\
\midrule
Electron & $4/3$ & $(3,\,2)$ & $4/3$ (fixed) \\
Muon & $16/5$ & $(5,\,4)$ & $3.2378\ldots$ \\
Tau & $25/9$ & $(9,\,5)$ & $25/9$ (fixed) \\
\bottomrule
\end{tabular}
\end{center}

The Koide-$r_\mu$ has no clean torus winding structure \SM.

\subsection*{Orthogonality of $\theta$ and $d/c$ \BM}

The $Z_3$ circulant has exactly two independent real parameters: the angle $\theta$
(rotation in the $Z_3$ space) and the ratio $d/c$ (eccentricity). Their orthogonality
is algebraically proved: $Q = \tfrac{1}{3} + \tfrac{d^2}{6c^2}$ does \emph{not}
contain $\theta$ --- the Koide condition is exclusively a statement about $d/c$,
independent of $\theta$ \BM.

Consequence for FFGFT: the corpus delivers $\theta = 2/9$ from the geometry of the
$Z_3$ circulant (Doc.~A110) \KM. The condition $d/c = \sqrt{2}$ is an \emph{additional,
independent} requirement --- the R113 bridge. Both together would be needed to derive
Koide fully from FFGFT. Currently only $\theta$ is anchored; $d/c = \sqrt{2}$ remains
open \SM.

\subsection*{Conclusions of the test}

\begin{enumerate}\itemsep4pt
\item \textbf{Angle $\theta = 2/9$ \KM:} The angular structure of the $Z_3$ circulant
  agrees between PDG, FFGFT and Doc.~A110 to $\lesssim 0.5\cdot\xi$.
\item \textbf{Amplitude $d/c$ --- two distinct objects:} \KM: given the observed
  generation-dependent $\varepsilon_i$, their effect on the Koide amplitude is
  calculable and brings the bare FFGFT $d/c$ ($+17\cdot\xi$ off $\sqrt{2}$)
  quantitatively to $\sqrt{2}$ --- 101\,\% explained (see §\,9).
  \SM: the FFGFT-native derivation of \emph{why} the $\varepsilon_i$ possess this
  generation-dependent structure --- and therefore why $d/c = \sqrt{2}$ should
  arise from the T$^4$/Z$_3$ geometry --- remains open (R113 bridge, §\,10).
\end{enumerate}

\section{Generation-Dependent Corrections and Koide}

\subsection*{The $\varepsilon_i$ are not uniform}

The bare FFGFT masses deviate from PDG masses by generation-dependent residuals $\varepsilon_i$:
\[
  m^{\rm PDG}_i = m^\circ_i \cdot (1 + \varepsilon_i)
\]

\begin{center}
\begin{tabular}{@{}lll@{}}
\toprule
Particle & $\varepsilon_i$ & Origin \\
\midrule
Electron & $+192\cdot\xi$ & SI projection, QED \\
Muon & $+152\cdot\xi$ & SI projection, QED extraction (R113) \\
Tau & $+73\cdot\xi$ & SI projection, measurement scatter \\
\bottomrule
\end{tabular}
\end{center}

The $\varepsilon_i$ decrease with generation --- they are not uniform.

\subsection*{Why non-uniform $\varepsilon_i$ shift Koide}

Koide is scale-invariant under uniform $\lambda$ \BM. For different $\varepsilon_i$,
to first order:
\begin{equation}
  \Delta Q \approx Q^\circ \cdot (\langle\varepsilon\rangle_m - \langle\varepsilon\rangle_{\sqrt{m}})
\end{equation}
Since $\varepsilon_\tau < \varepsilon_\mu < \varepsilon_e$ and $m_\tau$ dominates
the mass sum, $\langle\varepsilon\rangle_m < \langle\varepsilon\rangle_{\sqrt{m}}$,
hence $\Delta Q < 0$.

\subsection*{Numerical result}

\begin{center}
\begin{tabular}{@{}lll@{}}
\toprule
Quantity & Value & Status \\
\midrule
$Q({\rm bare}) - 2/3$ & $+8.06\cdot\xi$ & from quantum numbers \\
$\langle\varepsilon\rangle_m$ & $+77.7\cdot\xi$ & mass-weighted \\
$\langle\varepsilon\rangle_{\sqrt{m}}$ & $+90.0\cdot\xi$ & $\sqrt{m}$-weighted \\
$\Delta Q$ & $-8.23\cdot\xi$ & perturbation theory \\
$Q({\rm bare}) + \Delta Q - 2/3$ & $-0.17\cdot\xi$ & estimate \\
$Q({\rm PDG}) - 2/3$ & $-0.05\cdot\xi$ & exactly measured \\
Explanatory fraction & $101\,\%$ & \KM \\
\bottomrule
\end{tabular}
\end{center}

\subsection*{Revised conclusion \KM}

The statement ``ratios are correction-free'' holds for uniform corrections \BM.
For Koide additionally:
\begin{itemize}\itemsep2pt
\item Simple ratios $m_j/m_i$ are correction-free under uniform $\lambda$ \BM.
\item Koide $Q$ responds to non-uniform $\varepsilon_i$ because it is defined via
  $\sqrt{m_i}$.
\item The generation-dependent $\varepsilon_i$ (order $73$--$192\cdot\xi$, decreasing
  with generation) explain the transition $Q_{\rm bare}(+8\cdot\xi)\to Q_{\rm PDG}(-0.05\cdot\xi)$
  to $101\,\%$ \KM.
\item The remaining open point is: \emph{why do $\varepsilon_i$ decrease with generation?}
  The answer lies in the QED extraction (R113) and SI projection (Docs.~085, 333) ---
  their generation-dependent structure is analysed in §\,10.
\end{itemize}

\section{Why Do $\varepsilon_i$ Decrease with Generation?}
\label{sec:eps-origin}

\noindent\emph{Status separation in this section developed in correspondence with
Stefaan Vossen, applying the Constitutional Physics / Canonical Constitutional Record
framework (Doc.~351, 5--6~September~2026): the calculable effect of the observed
$\varepsilon_i$ \KM\ is distinguished from their source-native derivation \SM.}

\subsection*{Two sources of non-uniform corrections}

\paragraph{1.\ QED self-energy (R113) \QM.}
The measured pole mass of a charged lepton differs from the bare mass by the QED
self-energy. At leading order:
\begin{equation}
  \frac{\delta m_i}{m_i} \approx \frac{3\alpha}{4\pi}
  \ln\!\left(\frac{m_i}{\mu}\right) + \text{const}
\end{equation}
where $\mu$ is the renormalisation scale. Since $m_\tau \gg m_\mu \gg m_e$,
$\ln(m_i/\mu)$ grows with generation.

The observed trend $\varepsilon_e > \varepsilon_\mu > \varepsilon_\tau$ is,
however, \emph{opposite} to the naive QED self-energy that grows with mass.
The reason: $\varepsilon_i$ measures the gap bare FFGFT mass $\to$ PDG mass.
This gap is a \emph{difference} of two quantities each scaling with mass. The
SI projection contributes more strongly at small masses (relative correction
$\propto 1/m_i$ for a fixed absolute projection uncertainty).

\paragraph{2.\ SI projection (Docs.~085, 333) \SM.}
The SI projection maps the bare FFGFT mass $m^\circ_i = r_i\cdot\xi^{p_i}\cdot v\cdot K_{\rm frak}$
to the measured MeV value. The factor $K_{\rm frak} = 74/75$ is uniform \BM\ and
does not shift $Q$ (§\,5). The \emph{non-uniform} component arises from the
mass-dependent QED extraction: the experimental determination of $m_\mu$ from muon
decay and of $m_\tau$ from $\tau$ decay widths each involve different orders of
perturbation theory (R113). These are of order $\alpha\approx 1/137$, hence of
order $\xi$ to $100\xi$ --- consistent with the observed $\varepsilon_i$.

\subsection*{What FFGFT can derive and what it cannot}

\begin{center}
\begin{tabular}{@{}lll@{}}
\toprule
Quantity & Status & Origin \\
\midrule
Bare masses $m^\circ_i = r_i\cdot\xi^{p_i}\cdot v\cdot K_{\rm frak}$ & \KM & FFGFT internal \\
Uniform factor $K_{\rm frak}$ & \BM & FFGFT internal \\
Order of magnitude of $\varepsilon_i$ ($\sim\alpha/\pi$, order $\xi$--$100\xi$) & \KM & QED external \QM \\
Falling structure $\varepsilon_e > \varepsilon_\mu > \varepsilon_\tau$ & \SM &
  QED + SI projection, not derived from FFGFT quantum numbers \\
QED self-energy from T$^4$/Z$_3$ geometry & \SM & open bridge R113 \\
\bottomrule
\end{tabular}
\end{center}

\subsection*{Connection to Doc.~351}

The same non-uniform $\varepsilon_i$ structure explains the ratio deviations
in Doc.~351: $m_\mu/m_e$ deviates by $\sim 39\cdot\xi$ and $m_\tau/m_\mu$
by $\sim 78\cdot\xi$ from FFGFT predictions. Simple ratios are correction-free
under \emph{uniform} $\lambda$ \BM; under non-uniform $\varepsilon_i$ they carry
the difference residual $\varepsilon_j - \varepsilon_i$. The bridge R113 (QED
extraction, mass-dependent) is thus the common open point of Docs.~351 and~352
(Doc.~351, §\,5: ``R113: $m_\mu/m_e$ QED-dependent'', text correction
``model-independent'' $\to$ ``least model-dependent'').

\subsection*{Conclusion}

The observed falling $\varepsilon_i$ structure is consistent with a combined
QED/SI-projection account, but its derivation from FFGFT quantum numbers remains
open \SM. FFGFT delivers the bare masses; the spread of $\varepsilon_i$ is an
external \QM\ fact in its order of magnitude. Whether the QED self-energy is
derivable from the T$^4$/Z$_3$ geometry remains open \SM\ --- this is the R113 bridge.

\bigskip\noindent
\textbf{Verification script:}
\texttt{pruef\_352\_leptonreste.py} --- all numbers reproducible, pure Python standard library.

\section*{Sources}

PDG 2024 \QM. Doc.~006 (mass formula), Doc.~317 (lepton ladder values),
Doc.~351 (Galois identities, $\varepsilon_i$ values, R113), Doc.~A110
($Z_3$ circulant), Docs.~085, 333 (SI projection), R113 (QED dependence
$m_\mu/m_e$), Doc.~190/R114 (register entry for this document).
